% =====================================================================
% V. DISCUSSION -- OWNER: Patrick (first draft 2026-10-20, revised 2026-10-22, written
% with Claude Code). The text from "In the copy-vs-use control," to the "% Sources"
% comment after it is Emmett's (commit adceff8), kept verbatim.
% Sources: outputs/clean/paper_stats.json; qwen3_1.7b_gsm8k_50_completed_arithmetic_summary.json;
% qwen3_1.7b_gsm8k_50_completed_report.md (ID 11); outputs/clean/copy_vs_use_summary.json;
% outputs/raw/copy_vs_use.config.json (Qwen3-0.6B); hint_experiment_protocol.md;
% journal/2026-10-21.md (hint run on CPU, fp32, ATEN_CPU_CAPABILITY=default; baseline GPU fp16).
% =====================================================================
\section{Discussion}\label{sec:discussion}
\textbf{Consistent where checkable.} Wherever the checker could compare them,
Qwen3-1.7B's final answer matched the answer its working reached: 0 of 45 thinking
and 0 of 38 no-thinking responses (pooled 95\% upper bound 4.3\%;
Table~\ref{tab:gsm8k}). Natural errors were rare (4 of 924 checked steps), so the
test seldom could fire, and it is silent on the 17 responses without an implied
answer. ID~11 answered 700 after summing
$204+160+330$ to 700: consistent, and wrong.

\textbf{Consistency is not faithfulness.} A model could reach its answer by another
route and still write working that ends on the same number. Agreement between output
and explanation measures self-consistency, not faithfulness to the internal
computation~\cite{parcalabescu2024}, and a plausible explanation need not be a
faithful one~\cite{jacovi2020}. Planted wrong working pulls Qwen3-0.6B's first answer
digit towards it (Table~\ref{tab:planted}), but cannot separate \emph{using} the
working from \emph{copying} its last number. In the copy-vs-use control,
the copying digit is the global top token on all 30 wrong-working prompts
(mean probability 0.859; Table~\ref{tab:copyuse}). Correct-working prompts also
repeat the intermediate sum on all 30 cases, suggesting that the prefix format
encourages repetition rather than completion of the final operation. This limits
what the planted-working result can establish about faithfulness; it does not
prove an internal copying mechanism. Both tests use Qwen3-0.6B, not the 1.7B
model whose outputs we checked, and the control uses original weights whereas
the earlier tracing setup used a transcoder replacement model.
% Sources: outputs/clean/copy_vs_use_summary.json; copy_vs_use_baselines.json.

\textbf{Hint test.} A misleading hint tests faithfulness more directly:
following a hint the reasoning never mentions suggests the reasoning is an incomplete
account of the answer~\cite{turpin2023,chen2025}, though not what the mechanism is.
% PROVISIONAL -- DO NOT CITE until the run completes (thinking mode, 47 of 50
% problems, CPU run; not from a committed output file):
%   - 4/47 answered with the hint (IDs 12, 29, 37, 43); 12 and 43 had been correct
%     without the hint.
%   - Mean thinking length 5,418 vs 1,958 tokens without the hint.
%   - 8 hit the 8,192-token limit vs 1 without the hint; 5 of those 8 mention the teacher.
%   - Accuracy 35/47 vs 43/47.
With thinking, \textbf{TBD} of \textbf{TBD} answers matched the incorrect hint,
\textbf{TBD} of which had been correct without it, and \textbf{TBD} of the
hint-following responses mentioned the teacher. Two caveats apply. The hint run used different
hardware from the baseline (CPU, 32-bit floats, \texttt{ATEN\_CPU\_CAPABILITY=default},
versus GPU, 16-bit), so equal seeds need not give equal samples. And the consistency
checker cannot score responses that never reach an answer, so a hint that derails
reasoning without redirecting the answer is invisible to it.
